\documentclass[10pt,a4paper]{amsart}
\usepackage[margin=19mm]{geometry}
\usepackage{amsmath,amssymb,mathtools}
\usepackage[scr=rsfs,cal=euler]{mathalfa}
\usepackage{xfrac}
\usepackage[unicode,hidelinks,bookmarks=false]{hyperref}
\newcommand{\ie}{i.e.\ }
\newcommand{\cf}{cf.\ }
\newcommand{\resp}{respectively\ }
\newcommand{\Subsection}{\subsection}
\providecommand{\nobreakdash}{\nobreak\hbox{-}}
\setlength{\emergencystretch}{2em}
\multlinegap0pt
\usepackage{bm}
\usepackage{bbm}
\usepackage{dsfont}
\usepackage{xspace}
\usepackage{cancel}
\usepackage{mathtools}
\usepackage{amsthm}
\makeatletter
\@ifundefined{theorem}{\newtheorem{theorem}{Theorem}}{}
\@ifundefined{proposition}{\newtheorem{proposition}[theorem]{Proposition}}{}
\makeatother
\title[Analytic and Monte Carlo Approximations to the Distribution ]{Analytic and Monte Carlo Approximations to the Distribution of the First Passage Time of the Drifted Diffusion with Stochastic Resetting and Mixed Boundary Conditions}
\author{Juan Magalang, Riccardo Turin, Javier Aguilar, Laetitia Colombani, Daniel Sanchez-Taltavull, Riccardo Gatto}
\date{Source: arXiv:2311.03939v4 (CC BY 4.0)}
\begin{document}
\maketitle

\noindent\emph{Source and licence.} Analytic and Monte Carlo Approximations to the Distribution of the First Passage Time of the Drifted Diffusion with Stochastic Resetting and Mixed Boundary Conditions. Version arXiv:2311.03939v4 (CC BY 4.0), distributed under the Creative Commons Attribution 4.0 International licence, as recorded in the arXiv licence metadata for this exact version and shown on the abstract page https://arxiv.org/abs/2311.03939v4. Full rights notice: licenses/arxiv-2311.03939-CC-BY-4.0.txt.\par\smallskip

\noindent\textit{Editorial scope.} Counted unit: the Proposition whose source label is \texttt{prop:laplaceres} in the article (source0.tex lines 397--415, source label \texttt{prop:laplaceres}), delivered together with the article's own complete original proof of it (source0.tex lines 416--430, one \texttt{proof} environment of 15 lines). No mathematics is written, repaired or re-derived: statement and proof are the article's own text line for line. Required context retained uncounted: eq:sde\_with\_resets (lines 345--347); eq:fpt\_rv (lines 360--362); eq:kr (lines 1264--1267); eq:tildeSr (lines 1300--1303). Every definition, display and label of those lines is retained unchanged. Omitted from this package and not redistributed: 1--344, 348--359, 363--396, 431--1263, 1268--1299, 1304--1538. Nothing inside a retained excerpt is omitted or altered: every retained body is delivered exactly as the article's lines are written, so no omission receipt of any kind accompanies this package. Citations: the retained excerpts carry no citation command at all, so no bibliography entry of the article is transcribed and no reference is redistributed. Reference adaptation: none was needed and none was made; every reference inside the delivered unit resolves inside it, so no unresolved reference or citation is produced. Printed slips kept literal: no printed word, symbol, hypothesis or number of the counted statement or of its proof was altered. Layout and class: the article's own class is \texttt{elsarticle} with packages including amsmath, amsthm, amsfonts, amssymb, bbm, mathtools, appendix, dsfont, extarrows, enumerate, fullpage, xcolor, caption, subcaption; the wrapper is the lane's pinned \texttt{amsart} editorial wrapper at 10pt on A4, which restates the theorem-like environments the retained text needs and adds the article's own macro definitions that the retained text needs (none), with extra wrapper packages (bm, bbm, dsfont, xspace, cancel, mathtools). The delivered numbering is the unit's own. Assets: no retained span contains a figure environment, a graphics-inclusion command, a PDF-page inclusion, a table or a TikZ picture; no image file is copied into this package, the package declares no asset dependency and no image file is redistributed with it. Licence: version arXiv:2311.03939v4 of this article is distributed under the Creative Commons Attribution 4.0 International licence, as recorded in the arXiv licence metadata for this exact version and shown on the abstract page https://arxiv.org/abs/2311.03939v4. A full rights notice is delivered at licenses/arxiv-2311.03939-CC-BY-4.0.txt. Characters: every retained excerpt is pure ASCII, so the delivered engine, XeLaTeX with its default Latin Modern text font, reports no missing character.
\begin{equation} \label{eq:sde_with_resets}
	dX_t=(1-\chi_t) \cdot (v\,dt+\sqrt{2D}\,dW_t) +\chi_t \cdot (x_0-X_t),\quad \forall t\geq 0,
\end{equation}

\begin{equation}\label{eq:fpt_rv}
	\tau_r=\inf\{t\geq 0\, | \, X_t\leq0\}.
\end{equation}

\begin{equation}
\label{eq:kr}
	f_r(t)=-\partial_t{S}_r(t)\,,
\end{equation} 

\begin{equation} \label{eq:tildeSr}
	\tilde{S}_r(s)
	=\frac{1-\tilde{f}_0(s+r)}{s+r\tilde{f}_0(s+r)}\,,
\end{equation}

\begin{proposition}[Laplace transform of FPT distribution with resetting] \label{prop:laplaceres}
    Assume that $X$ solves the SDE with resetting of Eq. \eqref{eq:sde_with_resets}, with $X_0 = x_0 \in (0,1]$, $v \in \mathbb R$, $D>0$ and $r>0$. 
    Denote by $f_r$ the probability density of the absorption time $\tau_r$, defined in Eq. \eqref{eq:fpt_rv}. Then its Laplace transform is given by, $\forall s >-r$,
    \begin{align} \label{eq:tildef}
\tilde{f}_r(s) &
    = \frac{(s+r)\tilde{f}_0(s+r)}{s+r\tilde{f}_0(s+r)} \nonumber \\
    & 
    {\textstyle
    =\frac{(s+r) \left\{ \omega(s+r) \, \mathrm{cosh}[\theta(s+r)(x_0-1)]+v \, \mathrm{sinh}[\theta(s+r)(x_0-1)]\right\}}
    {s e^{\rho x_0}\left\{ \omega(s+r) \, \mathrm{cosh}[\theta(s+r)]-v \, \mathrm{sinh}[\theta(s+r)] \right\}
    + r \left\{ \omega(s+r) \, \mathrm{cosh}[\theta(s+r)(x_0-1)]+v \, \mathrm{sinh}[\theta(s+r)(x_0-1)] \right\} }},
\end{align}
where $\tilde{f}_0$ is given by, $\forall s> -v^2/(4D)$,
\begin{equation} \label{eq:tildef0}
	\tilde{f}_0(s)=
	e^{-\rho x_0}\frac{\omega(s) \, \mathrm{cosh}\{\theta(s)(x_0-1)\}+v \, \mathrm{sinh}\{\theta(s)(x_0-1)\}}
	{\omega(s) \, \mathrm{cosh}\{\theta(s)\}-v \, \mathrm{sinh}\{\theta(s)\}}\,.
\end{equation}
\end{proposition}

\begin{proof}
    We obtain from Eq.~\eqref{eq:kr} the Laplace transform of $f_r$ as
$\tilde{f}_r(s)=1-s\tilde{S}_r(s)$, for all $s \in \mathbb{R}$.
We combine it with Eq.~\eqref{eq:tildeSr} and it yields the first 
expression,
\begin{align*}
\forall s > -r, 	\; \tilde{f}_r(s)=\frac{(s+r)\tilde{f}_0(s+r)}{s+r\tilde{f}_0(s+r)}
 \; .
 \end{align*}
 This last expression holds $\forall s < -r$ under the assumption $ \tilde{f}_0(s+r)\neq -s/r$.  

We deduce the closed-form expression 
 in Eq.~\eqref{eq:tildef}
 by combining the first expression with $\tilde f_0$ in Eq.~\eqref{eq:tildef0}. Note that, for the function $\omega$ to be defined, we need $s+r \in \left(-v^2/(4D),\infty \right)$, viz. $s \in \left(- r -v^2/(4D), -r \right)$. 
\end{proof}

\end{document}
